% Source: https://artofproblemsolving.com/community/c3h2676017_2020_malaysia_imonst_1_primary_20_problems_25_hours_integer_0_answer_only
% Stable identity: aops-topic-2676017-q03
$P$ lies on the side of $BC$ away from the square, so $\angle ABP=90^\circ+60^\circ=150^\circ$. Since $BA=BC=BP$, triangle $ABP$ is isosceles with base angles $\dfrac{180^\circ-150^\circ}{2}=15^\circ$, so $\angle BPA=15^\circ$. Hence
\[\angle APC=\angle BPC-\angle BPA=60^\circ-15^\circ=\boxed{45^\circ}.\]

\begin{asy}
usepackage("amsmath");
size(3cm);
defaultpen(fontsize(10pt));
pair A=(0,0),B=(1,0),C=(1,1),D=(0,1),P=(1+sqrt(3)/2,0.5);
draw(A--B--C--D--cycle,linewidth(0.8));draw(B--P--C,linewidth(0.8));draw(A--P,dashed);
label("$A$",A,SW);label("$B$",B,S);label("$C$",C,N);label("$D$",D,NW);label("$P$",P,E);
void tk(pair a,pair b){pair m=(a+b)/2,u=0.06*unit(b-a);draw((m-rotate(90)*u)--(m+rotate(90)*u));}
tk(A,B);tk(B,C);tk(B,P);
\end{asy}
